\documentclass[10pt,a4paper]{amsart}
\usepackage[margin=19mm]{geometry}
\usepackage{amsmath,amssymb,mathtools}
\usepackage[scr=rsfs,cal=euler]{mathalfa}
\usepackage{xfrac}
\usepackage[unicode,hidelinks,bookmarks=false]{hyperref}
\newcommand{\ie}{i.e.\ }
\newcommand{\cf}{cf.\ }
\newcommand{\resp}{respectively\ }
\newcommand{\Subsection}{\subsection}
\providecommand{\nobreakdash}{\nobreak\hbox{-}}
\setlength{\emergencystretch}{2em}
\multlinegap0pt
\usepackage{bm}
\usepackage{bbm}
\usepackage{dsfont}
\usepackage{xspace}
\usepackage{cancel}
\usepackage{mathtools}
\usepackage{amsthm}
\makeatletter
\@ifundefined{theorem}{\newtheorem{theorem}{Theorem}}{}
\@ifundefined{proposition}{\newtheorem{proposition}[theorem]{Proposition}}{}
\makeatother
\def \B {\mathcal{B}}
\newcommand{\Ser}{\operatorname{Ser}}
\newcommand{\diame}{\mathsf{diam}}
\renewcommand{\leq}{\leqslant}
\title[Matroids, intersecting bases, and Borsuk property]{Matroids, intersecting bases, and Borsuk property}
\author{Gyivan L\'opez-Campos, Fr\'ed\'eric Meunier, Jorge L. Ram\'irez Alfons\'in}
\date{Source: arXiv:2506.17082v1 (CC BY 4.0)}
\begin{document}
\maketitle

\noindent\emph{Source and licence.} Matroids, intersecting bases, and Borsuk property. Version arXiv:2506.17082v1 (CC BY 4.0), distributed under the Creative Commons Attribution 4.0 International licence, as recorded in the arXiv licence metadata for this exact version and shown on the abstract page https://arxiv.org/abs/2506.17082v1. Full rights notice: licenses/arxiv-2506.17082-CC-BY-4.0.txt.\par\smallskip

\noindent\textit{Editorial scope.} Counted unit: the Proposition whose source label is \texttt{prop:ser} in the article (source0.tex lines 321--328, source label \texttt{prop:ser}), delivered together with the article's own complete original proof of it (source0.tex lines 330--345, one \texttt{proof} environment of 16 lines). No mathematics is written, repaired or re-derived: statement and proof are the article's own text line for line. Required context retained uncounted: none. Every definition, display and label of those lines is retained unchanged. Omitted from this package and not redistributed: 1--320, 329--329, 346--800. Nothing inside a retained excerpt is omitted or altered: every retained body is delivered exactly as the article's lines are written, so no omission receipt of any kind accompanies this package. Citations: the retained excerpts carry no citation command at all, so no bibliography entry of the article is transcribed and no reference is redistributed. Reference adaptation: none was needed and none was made; every reference inside the delivered unit resolves inside it, so no unresolved reference or citation is produced. Printed slips kept literal: no printed word, symbol, hypothesis or number of the counted statement or of its proof was altered. Layout and class: the article's own class is \texttt{amsart} with packages including amssymb, tikz, enumerate, enumitem, bbm, mathtools, hyperref, color, fullpage, subcaption, adjustbox, todonotes; the wrapper is the lane's pinned \texttt{amsart} editorial wrapper at 10pt on A4, which restates the theorem-like environments the retained text needs and adds the article's own macro definitions that the retained text needs (\texttt{\textbackslash B}, \texttt{\textbackslash Ser}, \texttt{\textbackslash diame}, \texttt{\textbackslash leq}), with extra wrapper packages (bm, bbm, dsfont, xspace, cancel, mathtools). The delivered numbering is the unit's own. Assets: no retained span contains a figure environment, a graphics-inclusion command, a PDF-page inclusion, a table or a TikZ picture; no image file is copied into this package, the package declares no asset dependency and no image file is redistributed with it. Licence: version arXiv:2506.17082v1 of this article is distributed under the Creative Commons Attribution 4.0 International licence, as recorded in the arXiv licence metadata for this exact version and shown on the abstract page https://arxiv.org/abs/2506.17082v1. A full rights notice is delivered at licenses/arxiv-2506.17082-CC-BY-4.0.txt. Characters: every retained excerpt is pure ASCII, so the delivered engine, XeLaTeX with its default Latin Modern text font, reports no missing character.
\begin{proposition}\label{prop:ser}
For every two matroids $M$ and $M'$, we have \[
f\bigl(\Ser(M,M')\bigr) \leq f(M) + f(M') \, .
\] Moreover, if $M$ has two disjoint bases that do not contain $p$ (the element used for the series connection), then
\[
f\bigl(\Ser(M,M')\bigr) \leq \min\bigl(f(M),f(M')\bigr) \, .
\]
\end{proposition}

\begin{proof}
Let $\B_1,\B_2,\ldots,\B_k$ be a partition of $\B(M)$ into $k=f(M)$ parts of diameter smaller than $\diame(\B(M))$, and $\B'_1,\B'_2,\ldots,\B'_{k'}$ be a partition of $\B(M')$ into $k'=f(M')$ parts of diameter smaller than $\diame(\B(M'))$. \\


[First inequality] We build a partition $\overline{\B}_1,\overline{\B}_2,\ldots,\overline{\B}_{k+k'}$ of $\B(\Ser(M,M'))$ as follows. Consider a basis $B \cup B'$ of $\Ser(M,M')$. If $B'$ does not contain $p$, then put $B \cup B'$ in $\overline{\B}_i$ with $i$ such that $B \in \B_i$. If $B'$ contains $p$, then put $B \cup B'$ in $\overline{\B}_{k+i}$ with $i$ such that $B' \in \B'_i$. We check that no part $\overline{\B}_i$ contains two bases of $\Ser(M,M')$ at distance $\diame\bigl(\B(\Ser(M,M'))\bigr)$ apart. Consider two bases $B_1 \cup B_1'$ and $B_2 \cup B_2'$ of $\Ser(M,M')$ that are at distance $\diame\bigl(\B(\Ser(M,M'))\bigr)$ apart. We distinguish three cases.

Case 1) $p \notin B_1' \cup B_2'$. Then $B_1$ and $B_2$ are at distance $\diame(\B(M))$ apart (since otherwise we could increase the distance between the two bases $B_1 \cup B_1'$ and $B_2 \cup B_2'$). Thus $B_1$ and $B_2$ do not belong to the same $\B_i$, and $B_1 \cup B_1'$ and $B_2 \cup B_2'$ do not belong to the same $\overline{\B}_i$.

Case 2) $p \in B_1' \triangle B_2'$. Without loss of generality, we have $p \in B_1' \setminus B_2'$. Then $B_1 \cup B_1'$ belongs to some $\overline{\B}_i$ with $i > k$, and $B_2 \cup B_2'$ belongs to some $\overline{\B}_i$ with $i \leq k$.

Case 3) $p \in B_1' \cap B_2'$. Then $B_1'$ and $B_2'$ are at distance $\diame(\B(M'))$ apart (since otherwise we could increase the distance between the two bases $B_1 \cup B_1'$ and $B_2 \cup B_2'$). Thus $B_1'$ and $B_2'$ do not belong to the same $\B'_i$, and $B_1 \cup B_1'$ and $B_2 \cup B_2'$ do not belong to the same $\overline{\B}_i$.\\

[Second inequality] We build a partition $\overline{\B}_1,\overline{\B}_2,\ldots,\overline{\B}_{\min(k,k')}$ of $\B(\Ser(M,M'))$ as follows.

Consider a basis $B \cup B'$ of $\Ser(M,M')$. If $k \leq k'$, put $B \cup B'$ in $\overline{\B}_i$ with $i$ such that $B \in \B_i$. Otherwise, put $B \cup B'$ in $\overline{\B}_i$ with $i$ such that $B' \in \B'_i$. We check that no part $\overline{\B}_i$ contains two bases of $\Ser(M,M')$ at distance $\diame\bigl(\B(\Ser(M,M'))\bigr)$ apart. Consider two bases $B_1 \cup B_1'$ and $B_2\cup B_2'$ of $\Ser(M,M')$ that are at distance $\diame\bigl(\B(\Ser(M,M'))\bigr)$ apart. Notice that if $M$ contains two disjoint bases that do not contain $p$, then $\diame\bigl(\B(\Ser(M,M'))\bigr) = \diame(\B(M)) + \diame(\B(M'))$. This implies that $B_1$ and $B_2$ are necessarily at distance $\diame(\B(M))$ apart and that $B_1'$ and $B_2'$ are necessarily at distance $\diame(\B(M'))$ apart. Therefore, the bases $B_1 \cup B_1'$ and $B_2 \cup B_2'$ cannot belong to the same $\overline{\B}_i$.
\end{proof}

\end{document}
